% Shared oracle page. Keep experiment-specific probe counts and results on page 3.
{\color{ink}\fontsize{19}{22}\selectfont\bfseries The stochastic gradient oracle\par}
{\small\color{muted}Reference page \enspace / \enspace random probes, unbiasedness and precision\par}
\vspace{7pt}

\textbf{The target is the gradient of the full GW objective.}
Let $r=P\one_m$, $c=P^\top\one_n$ and let $\odot2$ denote entrywise squares.
For symmetric $A,B$,
\[
 \begin{aligned}
 \nabla f(P)&=M(P)-4APB,\\
 M(P)&=2(A^{\odot2}r)\one_m^\top+2\one_n(B^{\odot2}c)^\top.
 \end{aligned}
\]
We compute $M(P)$ directly and estimate $APB$. The \emph{actual} marginals are
used: substituting $c=b$ at a column-infeasible iterate would give the wrong gradient.

\textbf{A probe is a random projection vector.}
For each oracle call, draw $q$ fresh, independent vectors
$z_s\in\{-1,+1\}^m$, with independent, equiprobable signs in every entry.
Stack them into $Z=[z_1,\ldots,z_q]\in\mathbb R^{m\times q}$.
Each vector involves all target points, and multiplication by $P$ and $A$ also
involves all source points. Thus $q$ counts projection vectors, not sampled cells
or coupling entries. The full coupling is retained with no imposed rank restriction.

\textbf{The implemented estimator.}
Using the same $Z$ in both terms, the oracle returns
\[
 \boxed{\widehat G(P;Z)=M(P)-\frac{2}{q}
 \left[A\bigl(P(BZ)\bigr)Z^\top+A(PZ)(BZ)^\top\right].}
\]
The computation forms $BZ$, $PZ$ and $P(BZ)$, multiplies the last two by $A$,
then forms the two $n\times m$ products above, avoiding explicit construction of
$APB$. Work grows with $q$; large sketches can cost more than the full gradient.

\textbf{Why there are two product terms.}
The estimator is the derivative of the following unbiased scalar sketch of GW:
\[
 \widehat f_Z(P)=r^\top A^{\odot2}r+c^\top B^{\odot2}c
 -\frac{2}{q}\left\langle PZ,\ A\bigl(P(BZ)\bigr)\right\rangle_F.
\]
Differentiating both occurrences of $P$ in the inner product produces the two
terms in $\widehat G$. The matrices $A,B$ and the objective being optimized remain
fixed. Reported losses are evaluations of the full $f$, rather than this random sketch.

\textbf{Conditional unbiasedness.}
The probes are independent of the current iterate and satisfy
$\mathbb E[ZZ^\top/q]=I_m$. Conditional on $P$,
\[
 \begin{aligned}
 \mathbb E\!\left[A(P(BZ))Z^\top/q\mid P\right]&=APB,\\
 \mathbb E\!\left[A(PZ)(BZ)^\top/q\mid P\right]&=APB,
 \end{aligned}
 \qquad
 \mathbb E[\widehat G(P;Z)\mid P]=\nabla f(P).
\]
The second identity uses $B=B^\top$. Fresh probes give the same unbiasedness
conditional on the entire optimization history before that draw.

\textbf{What the probe count controls.}
At a fixed $P$, the estimator averages $q$ independent single-probe estimates.
If $V(P)$ is the single-probe mean-squared Frobenius error, then
\[
 \mathbb E\!\left[\|\widehat G(P;Z)-\nabla f(P)\|_F^2\mid P\right]
 =\frac{V(P)}{q}.
\]
Four times as many probes halves the root-mean-square error at that same iterate.
This does not imply a corresponding improvement in loss or optimization time:
the iterates themselves change with $q$. Even $q\geq m$ does not make independent
random-sign probes exact, since $ZZ^\top/q$ need not equal $I_m$ for a finite draw.

\textbf{Connection to the update.}
FraMES averages outputs into $d_k$ and uses
\mbox{$H_k=d_k+\nabla(g_{\mu_k}\circ T)(P_k)$}. The penalty gradient is computed
without sampling; momentum is separate from the oracle.
\vfill
{\footnotesize\color{muted}Page 3 specifies the data, probe count, numerical precision and experimental results.}
